% ============================================================
% ch29.tex —— 第 29 章 方程动物园：笛卡尔的桥
% ============================================================
\chapter{方程动物园：笛卡尔的桥}

最后一篇，回到你最熟悉的场景：直角坐标系，抛物线 $y = x^2$。但你可能从没意识到，初中老师已经悄悄带你进了一扇大门——\textbf{你一直在给方程"拍照"}。这一章把这门手艺扶正：每个方程是一幅画，每幅画是一个方程。数学里最年轻的母亲（代数，四千年）与最古老的父亲（几何，更长）在此正式成婚，而证婚人是笛卡尔。

\section{动物园导游图}

1637 年，笛卡尔（和稍晚的费马，各自独立）发明坐标系时，脑子里想的是"让代数与几何互相打工"。四百年后的你在初中课堂上继承的正是这台翻译机：把方程的解描成点，点连成线。初中坐标系的存货，按"次数"重新布展：

\begin{center}
\begin{tabular}{clll}
\toprule
次数 & 方程 & 图像 & 备注 \\
\midrule
$1$ & $y = 2x + 1$ & 直线 & 斜率 $2$、截距 $1$ \\
$2$ & $y = x^2$ & 抛物线 & 开口向上 \\
$2$ & $x^2 + y^2 = 1$ & 单位圆 & 到原点距离 $1$ \\
$2$ & $xy = 1$ & 双曲线 & 两支，贴轴不碰轴 \\
$2$ & $x^2 + 2y^2 = 1$ & 椭圆 & 圆的拉伸版 \\
\bottomrule
\end{tabular}
\end{center}

\textbf{次数}（各项变量次数之和的最大值）是曲线的\textbf{辈分}：一次线辈分最低；二次的圆、抛物线、椭圆、双曲线是四堂兄弟（统称\textbf{圆锥曲线}——它们的统一点在选读小节揭晓）；三次起人口爆炸（本书主角椭圆曲线，第 32 章）。辈分为什么值钱？因为它决定\textbf{交点额度}——本章下半场与下一章的全部剧情。先看双曲线 $xy = 1$ 的一个执念：$x$ 越大 $y$ 越小，图像贴着 $x$ 轴和 $y$ 轴无限逼近却\textbf{永不触碰}（碰 $x$ 轴意味着 $y = 0$，即 $xy = 0 \ne 1$）。这两条被思念的直线叫\textbf{渐近线}——记住这份执念，第 31 章它要在"无穷远处"兑现。

\section{交点 $=$ 联立方程组的解}

\begin{kaiwei}
直线 $y = x$ 与圆 $x^2 + y^2 = 1$ 交于几点？直线 $y = x + 1$ 呢？$y = 1$ 呢？$y = 2$ 呢？
\end{kaiwei}

初中解法的标准姿势——代入消元。逐条走：

\paragraph{路一：$y = x$（过原点的对角线）。}代入圆：
\[
x^2 + x^2 = 1 \;\Longrightarrow\; 2x^2 = 1 \;\Longrightarrow\; x = \pm\frac1{\sqrt2}.
\]
两个交点 $\left(\tfrac1{\sqrt2}, \tfrac1{\sqrt2}\right)$ 与 $\left(-\tfrac1{\sqrt2}, -\tfrac1{\sqrt2}\right)$（对角线穿过圆心，交出一条直径）。判别式视角：消元后的二次方程 $\Delta = 0^2 - 4\cdot2\cdot(-1) = 8 > 0$，两实根。

\paragraph{路二：$y = x + 1$（斜切）。}代入：$x^2 + (x+1)^2 = 1$，展开合并：
\[
2x^2 + 2x = 0 \;\Longrightarrow\; 2x(x + 1) = 0 \;\Longrightarrow\; x = 0 \text{ 或 } x = -1 .
\]
两个交点 $(0, 1)$ 与 $(-1, 0)$（两点都在圆上 \checkmark）。$\Delta = 4 - 0 = 4 > 0$。

\paragraph{路三：$y = 1$（正切）。}代入：$x^2 + 1 = 1$，即 $x^2 = 0$——\textbf{重根}！$\Delta = 0$。几何上：直线在圆顶擦了一下，只有一个切点 $(0, 1)$。\textbf{判别式为零 $=$ 相切}——初中代数在这里说出了纯几何的语言。记住"切点算两票"（重根的代数体重），下一章它是主角。

\paragraph{路四：$y = 2$（离弦）。}代入：$x^2 + 4 = 1$，即 $x^2 = -3$，$\Delta = -12 < 0$——实数世界\textbf{无解}，图像上直线悬在圆上方，永不相遇。但慢着：$x^2 = -3$ 在\textbf{复数}世界有两个漂亮的解 $x = \pm\sqrt3\,\mathrm{i}$！也就是说，在"允许复数坐标"的扩大世界里，这条直线与圆\textbf{仍然交于两点}——$(\sqrt3\,\mathrm{i}, 2)$ 与 $(-\sqrt3\,\mathrm{i}, 2)$——只是交点住在实数画不出的楼层。

四条路合订成一张表：

\begin{center}
\begin{tabular}{ccccc}
\toprule
直线 & 消元后的 $\Delta$ & 实交点 & 复交点 & 几何姿势 \\
\midrule
$y = x$ & $8 > 0$ & $2$ & $2$ & 割（穿直径） \\
$y = x+1$ & $4 > 0$ & $2$ & $2$ & 割 \\
$y = 1$ & $0$ & $1$（重根） & $2$（计重） & 切 \\
$y = 2$ & $-12 < 0$ & $0$ & $2$ & 离 \\
\bottomrule
\end{tabular}
\end{center}

\section{第三条路的深意}

路四不是趣味题，它是下一章的命脉。看表最后一列：\textbf{复数世界里，交点数恒为 $2$，纹丝不动}。把结论钉好：

\begin{jielun}
直线与圆：实数交点 $\le 2$；\textbf{复数世界（计重数）恒等于 $2$}。
\end{jielun}

"恒等于"三个字价值连城：\textbf{交点不会凭空消失，只会搬家}——搬到复数看不见的楼层（路四），或在相切时叠成重票（路三）。代数的世界永远完整，是"实数画图"这个取景框太窄。第 11 章说过"原子清单随域而变"（$x^2+1$ 在实数上是原子、在复数上拆成两个），这里是它的几何翻版：\textbf{交点清单也随域而变}——而复数是那个"账目永平"的世界（代数基本定理：$n$ 次方程恰有 $n$ 个复根计重，第 11 章）。

\section{直线与抛物线：一样的账}

同一套流程在抛物线上再演一遍，验证规律不挑食：

$y = 0$ 与 $y = x^2$：$x^2 = 0$，重根——相切于原点，一票变两票（复数世界计 $2$）。

$y = -1$ 与 $y = x^2$：$x^2 = -1$，实数零点、复数两点（$\pm\mathrm{i}$）。

$y = x - 2$ 与 $y = x^2$：$x^2 - x + 2 = 0$，$\Delta = 1 - 8 = -7 < 0$——实数 $0$ 点、复数 $2$ 点。

\textbf{全套账目与直线--圆完全平行}：二次曲线 $\times$ 一次曲线 $=$ 恰好 $2$ 个复数交点（计重数）。而 $2 = 2\times1$——\textbf{次数在交点账本上是乘法}。那么自然的下一问：$2$ 次 $\times$ $2$ 次是不是恰好 $4$？$3$ 次 $\times$ $2$ 次是不是 $6$？这条"交点乘法表"能不能长成定理？——下一章，它将经历"实验支持、实验出卖、补洞成神"的三幕剧，最终登基为代数几何的第一条大定理（贝祖）。

\paragraph{选读小节：圆锥曲线的辈分秘密。}为什么圆、椭圆、抛物线、双曲线统称"圆锥曲线"？拿手电筒照墙面：光锥（圆锥）被墙面（平面）截出什么曲线，取决于切面角度——垂直照是圆，斜照是椭圆，切面与母线平行是抛物线，更陡穿过对顶锥是双曲线。两千年前阿波罗尼奥斯就是靠切圆锥给它们分的家。代数视角更利落：一般二次方程 $Ax^2 + Bxy + Cy^2 + Dx + Ey + F = 0$ 的判别式 $B^2 - 4AC$ 的符号（负、零、正）恰好分出椭圆、抛物线、双曲线三类——一个算术判据统一三兄弟。这个判别式与第 30、31 章的"曲线与无穷远的关系"直接挂钩：椭圆不碰无穷远、抛物线相切于无穷远、双曲线穿过无穷远——预告下一章的"画家平面"。

\begin{dongshou}
\begin{itemize}
  \yisi 对抛物线 $y = x^2$ 与直线 $y = mx + 1$ 逐个扫 $m$：联立得 $x^2 - mx - 1 = 0$，$\Delta = m^2 + 4 > 0$ 恒正——\textbf{任何斜率都割}！画图验证（这条直线族永远过 $(0,1)$，抛物线开口向上兜住它）。再改截距：$y = mx - 2$，$\Delta = m^2 - 8$，何时切？（$m = \pm2\sqrt2$。）
  \yier 求 $x^2 + y^2 = 1$ 与 $xy = 1$ 的实交点，并数复交点：代入 $y = 1/x$ 得 $x^2 + \tfrac{1}{x^2} = 1$，即 $x^4 - x^2 + 1 = 0$。对 $u = x^2$：$u^2 - u + 1 = 0$，$\Delta = -3 < 0$——实数零交点（图上两曲线确实互不搭理），但 $u = \tfrac{1\pm\sqrt3\,\mathrm{i}}2$，开方得四个复数 $x$，\textbf{复世界恰有 $2\times2 = 4$ 个交点}。账目与"次数相乘"的猜想再次吻合。
  \yier 双曲线 $xy = 1$ 与水平线 $y = k$ 交于几点？$k \ne 0$ 时一点（$x = 1/k$）；$k = 0$ 时零点。但代数上"每条直线该有 $2\times1 = 2$ 票"——横轴的两票去哪了？（第 31 章揭晓：横轴与双曲线两臂在\textbf{无穷远}处各有一票，合计两票，账目补平。先把这道"悬案"记在小本本上。）
\end{itemize}
\end{dongshou}

\begin{shaonao}
\textbf{给"相切 $=$ 重根"做个直接验证。}过 $y = x^2$ 上点 $(a, a^2)$ 的直线束 $y - a^2 = m(x - a)$。联立：
\[
x^2 - a^2 = m(x - a) \;\Longrightarrow\; (x-a)(x + a - m) = 0 .
\]
两个根：$x = a$（定根，直线本来就过这点）与 $x = m - a$（动根，随 $m$ 滑动）。\textbf{重合条件}：动根滑回定根，即 $m - a = a$，$m = 2a$——而 $2a$ 恰是 $y = x^2$ 在 $a$ 处的导数（切线斜率，第 15、16 章）！于是你用\textbf{纯代数}（不碰导数）重新推出了"切线斜率 $=$ 导数"：\textbf{切点是"两根重合"的极限位置}。这一手"根的追踪术"正是第 32 章椭圆曲线"弦切游戏"的入场券——书页之间，草蛇灰线。
\end{shaonao}

\begin{chengshi}
"直线与圆在复数世界恒交两点"需要"交点计重数"与"复数坐标点"两个约定，本章均已预演、未形式化；"二次曲线 $\times$ 一次曲线 $=$ $2$"对"直线落在曲线上"的退化情形（整条重合）不适用——贝祖定理（下一章）要明文排除"共享整条分支"。另外联立消元严格说要求"消元不丢失信息"（代入法对多项式方程是安全的，中学的"增根/失根"问题来自两边平方或约去含未知因式——本章操作均安全）。
\end{chengshi}

下一章，把"恰好 $mn$ 个"从口号变成定理——以及亲眼看它先被实验出卖（两个洞：切点与平行线），再被两个发明（重数与无穷远）救活，最终登基。
